% Ce que coûte la distance : latences mesurées entre trois régions simulées
% (outils/experiences/ch46/regions.py, délais réglés sur les médianes d'Azure).
\documentclass[border=6pt]{standalone}
\usepackage{sd}
\begin{document}
\begin{tikzpicture}[x=0.13mm, y=1mm, note/.append style={font=\sffamily\small},
    lab/.style={font=\sffamily\small, text=#1, align=left}, lab/.default=Dink]
  \foreach \y/\v/\t/\c in {
      0/0.1/{lecture sur la réplique locale (Singapour)}/Dteal,
     -7/2.9/{écriture à Paris, réplication asynchrone}/Dteal,
    -14/94.8/{écriture à Paris, synchrone vers la Virginie}/Dblue,
    -21/165.3/{écriture depuis Singapour sur le primaire de Paris}/Dblue,
    -28/168.8/{écriture à Paris, synchrone vers les deux répliques}/Dblue,
    -35/361.2/{nouvelle connexion Paris vers Virginie}/Dred,
    -42/657.8/{nouvelle connexion Singapour vers Paris}/Dred,
    -49/905.6/{nouvelle connexion Singapour vers Virginie}/Dred} {
    \node[note, anchor=east] at (-10,\y) {\t};
    \pgfmathsetmacro{\w}{max(\v,4)}
    \draw[draw=\c, fill=\c Bg, line width=0.5pt] (0,\y-2.2) rectangle (\w,\y+2.2);
    \node[note, anchor=west, text=\c] at (\w+8,\y) {\pgfmathprintnumber[use comma, fixed, fixed zerofill, precision=1, assume math mode=true]{\v}\,ms};}
  \foreach \t in {0,200,400,600,800,1000} {\draw[ligne] (\t,-53) -- (\t,-54.5);
    \node[note, below] at (\t,-54.5) {\t};}
  \draw[ligne] (0,-53) -- (1000,-53);
  \node[axe, anchor=north] at (500,-60) {millisecondes (médianes mesurées)};
  \foreach \t/\c in {88/Dslate, 162/Dslate} \draw[\c, dash pattern=on 2pt off 2pt, line width=0.4pt] (\t,3) -- (\t,-53);
  \node[lab=Dslate, anchor=south] at (88,3) {88};
  \node[lab=Dslate, anchor=south] at (162,3) {162};
  \node[lab=Dslate, anchor=south west] at (215,3) {allers-retours publiés par Azure (ms)};
\end{tikzpicture}
\end{document}
